% Standalone problem document, generated from the adjacent README.md.
% Compile with XeLaTeX. Mathematical content is maintained in Markdown.
\documentclass[11pt,a4paper]{article}
\usepackage[fontsize=10.5pt]{scrextend}
\usepackage[margin=22mm,headheight=15pt,headsep=9mm,footskip=11mm]{geometry}
\usepackage{fontspec}
\setmainfont{texgyrepagella}[Extension=.otf,UprightFont=*-regular,BoldFont=*-bold,ItalicFont=*-italic,BoldItalicFont=*-bolditalic]
\setsansfont{texgyreheros}[Extension=.otf,UprightFont=*-regular,BoldFont=*-bold,ItalicFont=*-italic,BoldItalicFont=*-bolditalic]
\setmonofont{texgyrecursor}[Extension=.otf,UprightFont=*-regular,BoldFont=*-bold,ItalicFont=*-italic,BoldItalicFont=*-bolditalic,Scale=MatchLowercase]
\usepackage{amsmath,amssymb}
\usepackage{unicode-math}
\setmathfont{texgyrepagella-math.otf}
\usepackage{xcolor,microtype,enumitem,needspace,fancyhdr}
\usepackage{bookmark}
\definecolor{ink}{HTML}{17344B}
\definecolor{accent}{HTML}{197D80}
\definecolor{muted}{HTML}{596773}
\hypersetup{colorlinks=true,linkcolor=accent,urlcolor=accent,pdftitle={MI-29 - Modulus-order
determinant comparison with an arbitrary base
power},pdfauthor={Open Problems in Numerical Linear Algebra}}
\urlstyle{same}
\setlength{\parindent}{0pt}
\setlength{\parskip}{5pt plus 1pt minus 1pt}
\linespread{1.04}
\setlength{\emergencystretch}{3em}
\widowpenalty=10000
\clubpenalty=10000
\displaywidowpenalty=10000
\allowdisplaybreaks[1]
\setlist{nosep,leftmargin=1.5em}
\providecommand{\tightlist}{\setlength{\itemsep}{0pt}\setlength{\parskip}{0pt}}
\providecommand{\pandocbounded}[1]{#1}
\pagestyle{fancy}
\fancyhf{}
\fancyhead[L]{\sffamily\footnotesize\color{muted}OPEN PROBLEMS IN NUMERICAL LINEAR ALGEBRA}
\fancyhead[R]{\sffamily\bfseries\footnotesize\color{accent}MI-29}
\fancyfoot[L]{\parbox[b]{0.9\textwidth}{\sffamily\fontsize{8}{10}\selectfont\color{muted}Editorial ratings; a bounded search cannot certify that a problem remains unsolved.\\Literature check: 2026-09-11}}
\fancyfoot[R]{\sffamily\footnotesize\color{muted}\thepage}
\renewcommand{\headrulewidth}{0.3pt}
\renewcommand{\headrule}{\hbox to\headwidth{\color{accent}\leaders\hrule height \headrulewidth\hfill}}
\makeatletter
\renewcommand{\section}{\Needspace{5\baselineskip}\@startsection{section}{1}{0pt}{12pt plus 2pt minus 2pt}{4pt}{\sffamily\bfseries\large\color{ink}}}
\renewcommand{\subsection}{\Needspace{4\baselineskip}\@startsection{subsection}{2}{0pt}{9pt plus 1pt minus 1pt}{3pt}{\sffamily\bfseries\normalsize\color{ink}}}
\makeatother
\setcounter{secnumdepth}{0}
\begin{document}
{\sffamily\small\color{muted}Matrix inequalities and norms\par}
\vspace{3pt}
{\raggedright\hyphenpenalty=10000\exhyphenpenalty=10000\sffamily\bfseries\fontsize{21}{24}\selectfont\color{ink}Modulus-order
determinant comparison with an arbitrary base power\par}
\vspace{7pt}
{\sffamily\small\textbf{Difficulty:} challenging\quad\textbf{Importance:} interesting
to specialist\par}
{\sffamily\small\bfseries\color{accent}Status: Solved\par}
\vspace{4pt}
{\color{accent}\hrule height 0.8pt}
\vspace{6pt}
\textbf{Rating rationale:} Arbitrary base powers and indefinite
Hermitian factors make extension of the squared-base theorem
challenging; the comparison has specialist importance for determinant
inequalities.

\subsection{Resolution --- 2026-09-11}\label{resolution-2026-09-11}

\textbf{Negative result by Matthew J. Colbrook}, Department of Applied
Mathematics and Theoretical Physics, University of Cambridge.
\textbf{Independent proof review: PASS.}

A rational positive definite \(A\) and invertible indefinite Hermitian
\(B\) in dimension three, with \(k=6\) and \(p=8\), reverse the proposed
determinant comparison. The exact right-minus-left gap is
\(21036678407451/156250000000000>0\). The known \(k=2\) theorem and the
variant \(B>0\) are not contradicted.

The exact target is resolved. The original statement and source evidence
are retained below; its former difficulty rating is historical.

\textbf{Primary manuscript:}
\href{https://github.com/ajt60gaibb/OpenProblemsInNLA/blob/main/matrix-inequalities-and-norms/MI-29/solution.pdf}{complete
proof PDF},
\href{https://github.com/ajt60gaibb/OpenProblemsInNLA/blob/main/matrix-inequalities-and-norms/MI-29/solution.tex}{standalone
TeX}, Theorem 1.1 and its proof;
\href{https://github.com/ajt60gaibb/OpenProblemsInNLA/blob/main/matrix-inequalities-and-norms/MI-29/solution.md}{authorship
and scope}. The
\href{https://github.com/ajt60gaibb/OpenProblemsInNLA/blob/main/references/colbrook-matrix-2026-09-11/verification/reviews/MI-29-review.md}{independent
review} checks the full original argument and records its hash. The
draft was AI-assisted; this is independent agent verification, not
external human peer review or formal certification.
\href{https://github.com/ajt60gaibb/OpenProblemsInNLA/blob/main/references/colbrook-matrix-2026-09-11/README.md}{Submission
record}.

\subsection{Problem statement}\label{problem-statement}

Does the inequality \[
\det(A^k+|AB|^p)\ge\det(A^k+|BA|^p)
\] hold for every integer \(n\ge1\), every positive definite
\(A\in\mathbb C^{n\times n}\), every invertible Hermitian
\(B\in\mathbb C^{n\times n}\), and every pair of real parameters
\(k,p\ge0\)? Here \(|X|=(X^*X)^{1/2}\), all powers of positive definite
matrices use spectral functional calculus, and \(X^0=I\).

Both determinants are positive real numbers. The same conjecture is
stated for semidefinite \(A\) and Hermitian \(B\) in the source;
invertible inputs avoid zeroth-power conventions, while the
positive-exponent singular cases follow by continuous approximation. No
commutation between \(A\) and \(B\) is assumed.

The two moduli have the same singular spectrum, but adding \(A^k\) tests
their alignment with the eigenvectors of the base matrix. The question
therefore concerns how spectral data and eigenvector geometry interact
in determinant estimates for matrix functions.

\newpage
\subsection{References}\label{references}

\begin{enumerate}
\def\labelenumi{\arabic{enumi}.}
\tightlist
\item
  M. M. Ghabries, \emph{Contributions to Matrix Inequalities and Some
  Applications}, PhD thesis, University of Angers and Lebanese
  University (2022), §2.5 and final Open Problems, Problem 2,
  pp.~109--110.
  \href{https://www.researchgate.net/publication/361793582_Contributions_to_Matrix_Inequalities_and_Some_Applications}{Thesis
  uploaded by its author}.
\item
  M. M. Ghabries, H. Abbas and B. Mourad, \emph{On some open questions
  concerning determinantal inequalities}, Linear Algebra Appl. 596
  (2020), 169--183, resolution of the earlier \(k=2\) question.
  \href{https://doi.org/10.1016/j.laa.2020.03.009}{Publisher}.
\item
  M. M. Ghabries, \emph{A log-majorization inequality for normal
  matrices with applications to determinantal inequalities and geometric
  means}, arXiv:2607.21163v1 (2026), §3.
  \href{https://arxiv.org/html/2607.21163}{Full text}.
\end{enumerate}

Status check (2026-09-10): the thesis states the full question after
proving \(k=2\) for all \(p\ge0\), and further cases when both matrices
are positive semidefinite. The July 2026 paper extends the known
squared-base result to normal matrices, but does not claim arbitrary
\(k\) here. Searches for the exact titles, ``modulus determinant
arbitrary power conjecture Ghabries'', and 2025/2026 found no complete
resolution. The negative-power companion in the thesis is not counted
separately.
\end{document}
